% !TeX root = ../教材题库_全部条目_转写工作稿.tex
% Source: XJ-B1-C06-CA-P227-Y1094
% 第6章 / 6.3 统计图表 / 教材 p.227、228、229、230 / PDF 第234、235、236、237页
% OCR未校对：公式与文字均可能有误；未完成数学TeX转写；禁止直接发给学生
% 原题及相邻原解截图（核对用）：../题目原图/XJ-B1-C06-CA-P227-Y1094-01.webp, ../题目原图/XJ-B1-C06-CA-P227-Y1094-02.webp, ../题目原图/XJ-B1-C06-CA-P227-Y1094-03.webp, ../题目原图/XJ-B1-C06-CA-P227-Y1094-04.webp
% 本文件只在转写完成后作为题干使用；OCR 条目尚不是可靠题干。
\begin{bankitem}{6.3 案例 案例3}{第6章 / 6.3 统计图表 / 教材 p.227、228、229、230 / PDF 第234、235、236、237页}{XJ-B1-C06-CA-P227-Y1094}{OCR未校对：公式与文字均可能有误；未完成数学TeX转写；禁止直接发给学生}
\begin{ocrdraft}
[教材p.227 / PDF第234页]
\par
案例3下面是某城市公共图书馆在一年中通过随机抽样调查得到的60天读者
\par
借书量（单位：册），并排序如下：
\par
213
\par
230
\par
239
\par
289
\par
291
\par
301
\par
308
\par
310
\par
311
\par
312
\par
318
\par
318
\par
337
\par
343
\par
344
\par
348
\par
349
\par
351
\par
360
\par
362
\par
368
\par
372
\par
374
\par
379
\par
383
\par
385
\par
390
\par
393
\par
396
\par
398
\par
425
\par
429
\par
438
\par
399
\par
400
\par
404
\par
406
\par
430
\par
436
\par
440
\par
441
\par
444
\par
446
\par
453
\par
456
\par
458
\par
471
\par
473
\par
475
\par
483
\par
484
\par
495
\par
498
\par
498
\par
521
\par
524
\par
549
\par
556
\par
568
\par
584
\par
为估计图书馆每天借书量的分布情况，以便合理安排工作人员，试根据以上数
\par
据制作一个频率分布表以帮助分析．
\par
（1）计算最大值与最小值的差（也称为这组数据的极差）
\par
[教材p.228 / PDF第235页]
\par
样本数据中最小值是213，最大值是584．它们的差是37L
\par
（2）确定组距和组数
\par
这60个数据散布在闭区间[213，584〕上．为了分组的方
\par
便，我们取一个略大的区间匚200，600），然后将该区间分成若
\par
一般来讲，当样本量
\par
干组．若取组距为50，那么
\par
为”时，可以参考经验公
\par
极差
\par
371
\par
式K=1+41g”将数据分
\par
组数一
\par
=7，42，
\par
组距
\par
50
\par
成K组．实际应用时，应
\par
因此可以将数据分为8组．
\par
当根据样本量的大小、数
\par
（3）将数据分组
\par
据的特点以及分析的要求
\par
将[200，600）八等分，所分8组为：
\par
来灵活确定．
\par
[200，250），[250，300），[300，
\par
350），
\par
[350，
\par
400），
\par
[400，450），匚450，500），[500，550），
\par
[550，
\par
600）．
\par
（4）列频率分布表
\par
当样本量是的观测数据中有刀，个落人第》组时，我们称一一是第组的
\par
频率计算出数据落人各组的频率为
\par
3
\par
2
\par
3
\par
3．3％，
\par
60
\par
60
\par
60
\par
列出频率分布表，如表6一3所示．
\par
表6一3
\par
分
\par
发生天数（频数玩
\par
频率\ensuremath{\text{[OCR:GBP]}}
\par
组
\par
[200，250）
\par
5％
\par
3
\par
[250，300）
\par
3．3％
\par
2
\par
[300，350）
\par
20％
\par
12
\par
[350，
\par
23．3％
\par
400）
\par
14
\par
[400，450）
\par
20％
\par
12
\par
[450，
\par
1\&3％
\par
500）
\par
11
\par
[500，550）
\par
5％
\par
3
\par
[550，600）
\par
3
\par
5％
\par
100％
\par
表6一3体现了样本数据落在各个小组的比例大小，从中可以看到，借书量在
\par
[350，400）内的天数最多，在[300，350）和[4佣，450）内的天数次之，大部分借书
\par
量集中在[300，500）之间．
\par
[教材p.229 / PDF第236页]
\par
数据的频率分布表初步展示了数据分布的一些规律．如果用图形来表示频率分
\par
布表就会更加形象和直观．例如，在平面直角坐标系中，用横轴表示读者借书量，
\par
纵轴表示频数，就可得到我们在初中学习过的频数分布直方图，如图6．3一4．
\par
频数
\par
12
\par
11
\par
3
\par
2
\par
0200250300350400450500550600讠卖者亻书量/\}
\par
图6．3一4
\par
如果我们在平面直角坐标系中，用横轴表示读者借书量，纵轴表示帮，将各
\par
分组的端点画在横轴上，用一组距
\par
作为小矩形的高，就得到由相连小矩形构成的
\par
图形．这样的图形称为频率分布直方图，如图6，3一5．
\par
组距
\par
0005
\par
0．004
\par
0\ensuremath{{}^\circ\mathrm C}03
\par
0，002
\par
0．001
\par
0200250300350400450500550600读者借书量/月牙
\par
图6．3一5
\par
频率
\par
图6．3一5中，每个小矩形的面积一组距\ensuremath{\times}
\par
一频率，所以各个小矩形的面积
\par
组距
\par
表示相应各组的频率，这样，频率分布直方图就以面积的形式反映了数据落在各个
\par
小组的频率的大小，并且容易知道，在频率分布直方图中各小矩形的面积之和等于l.
\par
从频率分布直方图可以直观地发现样本的一些分布规律，如在375附近达到
\par
“峰值”，并具有一定的对称性，这说明借书量在375册附近较为集中．另外还可以
\par
看出，特别少和特别多的借书量很少．
\par
要注意的是，同样一组数据，如果组距不同，得到的频率分布直方图的形状也
\par
[教材p.230 / PDF第237页]
\par
会不同．
\par
如果将频率分布直方图中的左边和右边各延长一个分组，取各相邻小矩形上底
\par
边的中点，用线段顺次连接各点，就得到频率分布折线图．频率分布折线图也反映
\par
出数据频率分布的规律．图6，3一6是图书馆读者借书量的频率分布折线图．
\par
组距
\par
0005
\par
0．004
\par
0003
\par
0．002
\par
0001
\par
200250300350400450500550600650i卖者亻昔《量/川\}
\par
图6．3一6
\par
统计图表是展示数据的有效工具，通过制作统计图表可使杂乱的数据有条理地、
\par
形象地显示出来，这将为研究实际问题、做好统计分析打下良好的基础．
\end{ocrdraft}
\end{bankitem}
